% 原书第 203—206 页；第 13.3 节全文。
\section{有匹配随机试验中的靶向估计}
\label{sec:ch13-with-matching}

回顾假设的社区随机试验所采用的配对匹配方案。首先，选出可能纳入研究的单位。随后，对候选单位的基线协变量应用匹配算法，形成 $15$ 个最佳匹配对。在这些匹配对内随机分配干预，再通过纵向随访测量结局。这种配对匹配方案被认为是自适应的，因为所得匹配对是所有候选单位的基线协变量的函数（van der Laan et al.，2013a；Balzer et al.，2015，2016c）。这种设计也称为“非二分匹配”（nonbipartite matching）和“最优多变量匹配”（optimal multivariate matching；Greevy et al.，2004；Zhang and Small，2009；Lu et al.，2011）。

这种自适应设计使数据产生相依性。由于匹配对的构造是全部 $n$ 个研究单位的基线协变量的函数，观测数据并不是 $n/2$ 个独立同分布的成对观测；而当前实践有时作了这样的假设，例如 Klar and Donner（1997）、Freedman et al.（1997）、Campbell et al.（2007）以及 Hayes and Moulton（2009）。实际上，我们有 $O=(W,A,Y)$ 的 $n$ 个相依副本。尽管如此，数据中仍保留了大量条件独立性。主要而言，只要将研究单位的基线协变量视为固定，就重新得到 $n/2$ 个条件独立的单位：
\[
 \bar O_j=(O_{j1},O_{j2})
 =((W_{j1},A_{j1},Y_{j1}),(W_{j2},A_{j2},Y_{j2})),
\]
其中，指标 $j=1,\ldots,n/2$ 表示按照基线协变量 $(W_1,\ldots,W_n)$ 的相似性，将候选单位 $\{1,\ldots,n\}$ 划分为匹配对。下文始终用下标 $j1$ 和 $j2$ 表示匹配对 $j$ 内的观测。给定研究单位的基线协变量，观测数据的条件分布分解为
\[
 \begin{aligned}
 &P_0(O_1,\ldots,O_n\mid W_1,\ldots,W_n)\\
 &\quad=\prod_{j=1}^{n/2}
   P_0(A_{j1},A_{j2}\mid W_1,\ldots,W_n)
   \times P_0(Y_{j1}\mid A_{j1},W_{j1})\\
 &\qquad\qquad\qquad{}\times P_0(Y_{j2}\mid A_{j2},W_{j2})\\
 &\quad=\prod_{j=1}^{n/2}0.5
   \times P_0(Y_{j1}\mid A_{j1},W_{j1})
   \times P_0(Y_{j2}\mid A_{j2},W_{j2}),
 \end{aligned}
\]
第二个等式来自匹配对内的干预随机分配。对于总体效应（PATE）的估计与推断，需要假设每个社区的基线协变量 $W_i$ 都独立抽自某个共同分布 $P_0(W)$。对于样本效应（SATE）的估计与推断，则可以放宽这一协变量分布假设（Balzer et al.，2016c）。

尽管数据存在相依性，总体或样本效应的 TMLE 仍可忽略配对匹配设计来实施（van der Laan et al.，2013a；Balzer et al.，2016c）。换言之，按照第 13.2 节概述的过程就可以得到点估计。在步骤 1，使用事先指定的参数工作模型，或者下文详细介绍的更具数据自适应性的方法，得到条件结局均值的初始估计量。在步骤 2，利用已知或估计的暴露机制中的信息，对初始估计量进行靶向更新。最后，在步骤 3，预测所有观测分别在处理与对照条件下的结局，再对这些靶向预测值之差取样本平均。

在采用自适应配对匹配的试验中，TMLE 同时是总体效应和样本效应的渐近正态估计量（van der Laan et al.，2013a；Balzer et al.，2016c）。对于 PATE，可以用无匹配试验中估计影响曲线的样本方差
$\frac1n\sum_{i=1}^n[D_n^{\mathcal P}(g_n,\bar Q_n^*)(O_i)]^2$
除以 $n$，来估计其方差。不过，这个方差估计量忽略了配对匹配所带来的任何精度改善，在合理假设下会偏于保守。如果考虑匹配对内残差可能存在的相关性，则可得到不那么保守的方差估计量：
\begin{equation}
 \begin{aligned}
 \rho_n(\bar Q_n^*)(\bar O_j)
 =\frac1{n/2}\sum_{j=1}^{n/2}
  &(Y_{j1}-\bar Q_n^*(A_{j1},W_{j1}))\\
  &{}\times(Y_{j2}-\bar Q_n^*(A_{j2},W_{j2})).
 \end{aligned}
 \tag{13.7}
\end{equation}
这一结果见 van der Laan et al.（2013a）。在有限样本中，对于配对匹配下总体效应的 TMLE，我们建议用下式估计方差：
\[
 \bar\sigma_n^{2,\mathcal P}
 =\frac{\dfrac1n\displaystyle\sum_{i=1}^n
           [D_n^{\mathcal P}(g_n,\bar Q_n^*)(O_i)]^2
          -2\rho_n(\bar Q_n^*)(\bar O_j)}{n}.
\]

在配对匹配试验中，TMLE 减去样本效应（SATE），表现为一个影响曲线的经验均值，而该影响曲线依赖于不可识别的量（Balzer et al.，2016c）。尽管如此，其影响曲线的一个保守的代入估计量为
\[
 \bar D_n^S(g_n,\bar Q_n^*)(\bar O_j)
 =\frac12\left[D_n^S(g_n,\bar Q_n^*)(O_{j1})
              +D_n^S(g_n,\bar Q_n^*)(O_{j2})\right],
\]
其中，$D_n^S(g_n,\bar Q_n^*)(O)$ 为无匹配试验中观测 $O$ 的估计影响曲线，见式 (13.4)。在有限样本中，用所估计的成对影响曲线的样本方差除以 $n/2$，来保守估计样本效应的 TMLE 的方差：
\[
 \bar\sigma_n^{2,S}
 =\frac{\dfrac1{n/2}\displaystyle\sum_{j=1}^{n/2}
       [\bar D_n^S(g_n,\bar Q_n^*)(\bar O_j)]^2}{n/2}.
\]

\noindent\textbf{步骤 1 的预先指定自适应方法：初始估计。}
配对匹配通过平衡各干预组中决定结局的基线因素，提高研究效率，例如 Imai et al.（2009）、van der Laan et al.（2013a）以及 Balzer et al.（2015，2016c）。尽管如此，基线预测因素往往仍有残余不平衡，在分析中调整这些协变量，可以进一步提高效率。在贯穿本章的例子中，假设在测量基线 HIV 流行率之前，就已经形成了匹配对。因此，基线流行率很可能在不同匹配对间存在变异，而它是 HIV 发生率的一个已知驱动因素。在分析中调整基线流行率，很可能通过两个机制提高功效：（1）减小用于点估计的 TMLE 的方差；（2）使方差估计量不那么保守。不幸的是，事先并不清楚，调整流行率是否比调整其他协变量，例如男性包皮环切覆盖率或社区层面的 HIV RNA 病毒载量指标，能够带来更高的功效。只有 $n/2=15$ 个条件独立单位，限制了调整集的大小。调整过多协变量会导致过拟合。与前面一样，我们希望以数据自适应方式，选择使经验效率最大的候选 TMLE，即工作回归模型。

对于无匹配试验，初始估计条件结局均值 $\bar Q_0(A,W)$ 的数据自适应方法，可以修改后用于配对匹配试验。如前所述，需要预先指定候选估计量库、表现度量以及交叉验证方案。初始估计条件结局均值 $\bar Q_0(A,W)$ 时，可以使用相同的候选工作模型库。不过，为衡量表现，希望以配对匹配下 TMLE 的估计方差作为风险。为进一步说明，考虑无匹配试验中样本效应的损失函数。最小化残差平方和，即式 (13.6)，是以条件结局均值 $\bar Q_0(A,W)$ 为目标。因此，算法可能选择一个调整某个协变量的工作模型；该协变量虽然能很好地预测结局，但我们在设计中已对其实现了完全匹配。在贯穿本章的例子中，假设在地区内对社区配对，因为预期 HIV 发生率在地区间高度异质。因此，最小化 $D_n^S(g_0,\bar Q)$ 的经验方差，可能导致选择包含干预和地区主效应的候选 TMLE。相对于未经调整的算法，这样的选择不会提高分析精度，因为我们已经在设计中“控制”了地区。我们真正希望选择的是，在配对匹配试验中，使感兴趣参数的估计精度最大的候选 TMLE。因此，PATE 的损失函数为
\begin{equation}
 \begin{aligned}
 \bar{\mathcal L}^{\mathcal P}(g_0,\bar Q)(\bar O_j)
 ={}&\frac12\{D_n^{\mathcal P}(g_0,\bar Q)(O_{j1})\}^2
    +\frac12\{D_n^{\mathcal P}(g_0,\bar Q)(O_{j2})\}^2\\
    &{}-2(Y_{j1}-\bar Q(A_{j1},W_{j1}))
          (Y_{j2}-\bar Q(A_{j2},W_{j2})),
 \end{aligned}
 \tag{13.8}
\end{equation}
SATE 的损失函数为
\begin{equation}
 \bar{\mathcal L}^S(g_0,\bar Q)(\bar O_j)
 =\{\bar D_n^S(g_0,\bar Q_n^*)(\bar O_j)\}^2.
 \tag{13.9}
\end{equation}
同样，我们将暴露机制视为已知：$g_0(A\mid W)=0.5$。

最后，在交叉验证方案中，应将匹配对视为条件独立单位。换言之，将数据划分为 $V$ 折时，应保留配对关系。在小型试验中，留一对交叉验证往往是合适的。经过这些修改，就可以实施第 13.2 节概述的交叉验证方案，以数据自适应方式选择使配对匹配试验中 TMLE 的估计方差最小的候选工作模型。随后，同样用全部数据重新拟合所选工作模型，并用它预测所有观测分别在处理和对照条件下的结局。对预测结局之差取平均，就给出干预效应的估计。

